\section{Threat-Transport Graph and Hinge-Zone Diagnosis}
Etosha should not be read as a flat field of independent hotspots. It behaves as a transport network in which illegal pressure enters from a limited set of ingress opportunities, moves through a sparse relay structure, and ultimately threatens high-value habitat sinks. The point of the diagnostic layer is therefore to infer where weighted loss can travel, not merely where static ecological value is high.

\subsection{Network Abstraction}
We partition the park into an \(8\times 8\) management grid and construct two linked graphs:
\begin{itemize}
  \item a \textbf{road and access graph} \(\mathcal{G}_R=(V_R,E_R,w)\) for station travel and patrol mobility;
  \item a \textbf{transport graph} \(\mathcal{G}_T=(V_T,E_T)\) whose nodes represent management cells, ingress cells, and habitat sinks.
\end{itemize}

Each management cell \(i\) stores four strategic attributes:
\[
V_i:\text{ biological consequence},\quad
A_i:\text{ ingress accessibility},\quad
C_i:\text{ transport criticality},\quad
S_i:\text{ service slack}.
\]
All map plates in this section are anchored to the official Etosha basemap and upgraded with GIS overlays so that the diagnostic outputs stay visibly tied to real park geography.

\subsection{Ingress Prior}
The first question is not which cell is valuable, but where pressure is most likely to enter. Let \(\mathcal{E}\) denote the candidate ingress set. For entry node \(e\in\mathcal{E}\), define a prior weight
\begin{equation}
\pi_e=
\frac{\exp(\theta_1 g_e+\theta_2 \ell_e-\theta_3 p_e)}
{\sum_{k\in\mathcal{E}}\exp(\theta_1 g_k+\theta_2 \ell_k-\theta_3 p_k)},
\label{eq:prior}
\end{equation}
where \(g_e\) measures access convenience, \(\ell_e\) measures downstream biological draw, and \(p_e\) measures recent defensive pressure near the entry corridor. Equation~\eqref{eq:prior} attaches uncertainty to entry structure first and lets the network propagate that pressure inland.

\subsection{Transport Criticality}
Let \(Q\) be the transient transition matrix on the non-sink portion of \(\mathcal{G}_T\). The fundamental matrix
\[
N=(I-Q)^{-1}
\]
gives the expected number of visits to each transient node before absorption at a habitat sink. We define the transport criticality of cell \(i\) by
\begin{equation}
C_i = \frac{\sum_{e\in\mathcal{E}} \pi_e N_{ei}-\min_k \sum_{e\in\mathcal{E}} \pi_e N_{ek}}
{\max_k \sum_{e\in\mathcal{E}} \pi_e N_{ek}-\min_k \sum_{e\in\mathcal{E}} \pi_e N_{ek}+\varepsilon}.
\label{eq:criticality}
\end{equation}
Cells with high \(C_i\) are not necessarily those with the highest wildlife value. They are the cells through which an outsized share of weighted threat flow is expected to pass before it can be neutralized.

\subsection{Service Slack}
Response quality is measured here as slack relative to a service commitment rather than as raw travel time alone. Let \(v(i)\in V_R\) be the access node for management cell \(i\), and let \(\mathcal{S}\) be the station set. The shortest-path distance is
\[
d_i=\min_{s\in\mathcal{S}} \operatorname{dist}_{\mathcal{G}_R}(s,v(i)).
\]
If the dispatch-delay term is approximated by a normal perturbation with mean \(\mu_\delta\) and standard deviation \(\sigma_\delta\), then the probability that service arrives before the commitment threshold \(\tau^\star\) is
\begin{equation}
S_i = \Phi\!\left(\frac{\tau^\star-d_i-\mu_\delta}{\sigma_\delta}\right),
\label{eq:response}
\end{equation}
where \(\Phi(\cdot)\) is the standard normal distribution function. High \(S_i\) means a zone still has operational slack; low \(S_i\) means the zone is structurally late even before randomness is added.

\begin{figure}[H]
  \centering
  \includegraphics[width=\linewidth]{fig02_etosha_grid.png}
  \caption{Etosha transport abstraction on the official basemap, showing management cells, station anchors, and the geographic skeleton used for ingress-to-habitat analysis.}
  \label{fig:grid}
\end{figure}

\subsection{Biological Consequence}
We measure wildlife consequence independently of network structure:
\begin{equation}
V_i = 0.45W_i + 0.35H_i + 0.20S_i^{(\mathrm{bio})},
\label{eq:value}
\end{equation}
where \(W_i\) is wildlife density, \(H_i\) is water dependence, and \(S_i^{(\mathrm{bio})}\) is species sensitivity. This term answers a different question from transport criticality. It asks how expensive a loss would be if it reached cell \(i\), not how likely the path is to pass there.

\subsection{Residual Defense Debt}
The main diagnostic output is not a standard risk score but a residual defense debt:
\begin{equation}
\Gamma_i = V_i\Bigl(0.50C_i + 0.30A_i + 0.20(1-S_i)\Bigr).
\label{eq:risk}
\end{equation}
\(\Gamma_i\) becomes large when a cell is biologically consequential, strategically central to loss flow, easy to reach from ingress structure, and operationally short on service slack. This is the quantity the controller will later try to suppress.

\subsection{Etosha Diagnosis Under the New Framework}
The resulting picture is materially different from a simple corridor heatmap. The east and south still matter, but not because every cell in those regions is equally dangerous. What matters is a smaller hinge system: a set of eastern ingress-facing cells, a southern relay strip, and several waterhole-linked sink connectors. These are the cells where loss flow can still remain both biologically expensive and operationally under-served.

The practical implication is sharp. Etosha should not be defended by uniform corridor saturation. It should be defended by breaking the continuity of that hinge system.

\begin{figure}[H]
  \centering
  \includegraphics[width=\linewidth]{fig03_risk_heatmap.png}
  \caption{Residual defense-debt field on the official basemap, emphasizing hinge zones and transport bottlenecks rather than broad undifferentiated hotspots.}
  \label{fig:riskmap}
\end{figure}

\begin{figure}[H]
  \centering
  \includegraphics[width=\linewidth]{fig04_risk_decomposition.png}
  \caption{Decomposition of Etosha's loss-flow debt into biological consequence, transport criticality, ingress accessibility, and service slack.}
  \label{fig:riskdecomp}
\end{figure}
\FloatBarrier
